% Part: incompleteness
% Chapter: arithmetization-syntax
% Section: proofs-in-nd

\documentclass[../../../include/open-logic-section]{subfiles}

\begin{document}

\olfileid{inc}{art}{pnd}
\olsection{\usetoken{P}{derivation} در استنتاج طبیعی}

\begin{explain}
برای حسابی‌سازی !!{derivation}s باید !!{derivation}s را به‌صورت اعداد بازنمایی
کنیم. چون !!{derivation}s درخت‌هایی از !!{formula}s هستند که هر استنتاج
یک یا دو برچسب دارد، بازنمایی بازگشتی آشکارترین رویکرد است:
!!{derivation} را با چندتایی‌ای بازنمایی می‌کنیم که مؤلفه‌هایش تعداد
زیر!!{derivation}sی بی‌واسطهٔ منتهی به مقدمات آخرین استنتاج،
بازنمایی‌های این زیر!!{derivation}s، !!{formula} پایانی، برچسب رفعِ آخرین
استنتاج و عددی نشان‌دهندهٔ نوع آخرین استنتاج‌اند.
\end{explain}

\begin{defn}
  اگر $\delta$ !!a{derivation}ی در استنتاج طبیعی باشد، $\Gn{\delta}$ به
  طور استقرایی چنین تعریف می‌شود:
  \begin{enumerate}
  \item
    اگر $\delta$ فقط از فرض~$!A$ تشکیل شده باشد، $\Gn{\delta}$ برابر
    $\tuple{0,\Gn{!A},n}$ است. عدد~$n$ برای فرض !!{undischarged} برابر
    $0$، و در غیر این صورت برچسب عددی آن است.
  \item
    اگر $\delta$ با استنتاجی دارای صفر، یک، دو یا سه مقدمه پایان یابد،
    آنگاه $\Gn{\delta}$ به‌ترتیب برابر است با
    \begin{align*}
      & \tuple{0, \Gn{!A}, n, k}, \\
      & \tuple{1, \Gn{\delta_1}, \Gn{!A}, n, k}, \\
      & \tuple{2, \Gn{\delta_1}, \Gn{\delta_2}, \Gn{!A}, n, k}, \text{ یا}\\
      & \tuple{3, \Gn{\delta_1}, \Gn{\delta_2}, \Gn{\delta_3}, \Gn{!A}, n,
        k}.
    \end{align*}
    اینجا $\delta_1$، $\delta_2$ و~$\delta_3$ زیر!!{derivation}sیی هستند
    که به مقدمه یا مقدمات آخرین استنتاج در~$\delta$ منتهی می‌شوند؛
    $!A$ نتیجهٔ آخرین استنتاج در~$\delta$ است؛ $n$ برچسب رفعِ آخرین
    استنتاج است ($0$ اگر استنتاج هیچ فرضی را رفع نکند)؛ و~$k$ بر حسب
    قاعدهٔ آخرین استنتاج از جدول زیر به دست می‌آید.

    \begin{tabular}{lccccccc}
      \text{قاعده:} & \Intro{\land} & \Elim{\land} & \Intro{\lor} & \Elim{\lor} \\
      $k$: & 1 & 2 & 3 & 4  \\[1.5ex]
      \text{قاعده:} &  \Intro{\lif} & \Elim{\lif} & \Intro{\lnot} & \Elim{\lnot} \\
      $k$: & 5 & 6 & 7 & 8 \\[1.5ex]
      \text{قاعده:}  &  \FalseInt & \FalseCl & \Intro{\lforall} & \Elim{\lforall} \\
      $k$: & 9 & 10 & 11 & 12 \\[1.5ex]
      \text{قاعده:} & \Intro{\lexists} & \Elim{\lexists} & \Intro{\eq} & \Elim{\eq} \\
      $k$: & 13 & 14 & 15 & 16
    \end{tabular}
  \end{enumerate}
\end{defn}

\begin{ex}
  !!{derivation} بسیار سادهٔ زیر را در نظر بگیرید:
  \begin{prooftree}
    \AxiomC{$\Discharge{!A \land !B}{1}$}
    \RightLabel{\Elim{\land}}
    \UnaryInfC{$!A$}
    \DischargeRule{\Intro{\lif}}{1}
    \UnaryInfC{$(!A \land !B) \lif !A$}
  \end{prooftree}
  عدد گودل فرض برابر
  $d_0=\tuple{0,\Gn{!A\land !B},1}$ خواهد بود. عدد گودل
  !!{derivation} منتهی به نتیجهٔ~$\Elim{\land}$ برابر
  $d_1=\tuple{1,d_0,\Gn{!A},0,2}$ خواهد بود ($1$ چون
  $\Elim{\land}$ یک مقدمه دارد، سپس عدد گودل نتیجهٔ~$!A$، $0$ چون
  هیچ فرضی رفع نمی‌شود، و~$2$ عدد رمزکنندهٔ~$\Elim{\land}$ است).
  بنابراین عدد گودل کل !!{derivation} برابر
  $\tuple{1,d_1,\Gn{((!A\land !B)\lif !A)},1,5}$ است؛ یعنی
  \[
  \tuple{1, \tuple{1, \tuple{0, \Gn{(!A \land !B)}, 1}, \Gn{!A}, 0, 2},
    \Gn{((!A \land !B) \lif !A)}, 1, 5}.
  \]
\end{ex}

\begin{explain}
پس از تثبیت بازنمایی !!{derivation}s باید همچنین نشان دهیم می‌توانیم
اعداد گودل چنین !!{derivation}sیی را به‌طور بازگشتی اولیه دست‌کاری و
ویژگی‌ها و روابط اساسی‌شان را بیان کنیم. بعضی عملیات ساده‌اند: برای
نمونه، با داشتن عدد گودل~$d$ از !!a{derivation}ی،
$\fn{EndFmla}(d)=(d)_{(d)_0+1}$ عدد گودل !!{formula} پایانی،
$\fn{DischargeLabel}(d)=(d)_{(d)_0+2}$ برچسب رفع، و
$\fn{LastRule}(d)=(d)_{(d)_0+3}$ عدد نشان‌دهندهٔ نوع آخرین استنتاج را
به دست می‌دهد. بعضی موارد بسیار دشوارترند و دست‌کم طرحی از چگونگی
انجامشان می‌دهیم. هدف نشان‌دادن این است که رابطهٔ «$\delta$
!!a{derivation}ی از~$!A$ بر پایهٔ~$\Gamma$ است» رابطه‌ای بازگشتی اولیه
میان اعداد گودل~$\delta$ و~$!A$ است.
\end{explain}

\begin{prop}
  روابط زیر بازگشتی اولیه‌اند:
  \begin{enumerate}
  \item $!A$ با برچسب~$n$ به‌عنوان فرض در~$\delta$ رخ می‌دهد.
  \item همهٔ فرض‌های~$\delta$ با برچسب~$n$ به صورت~$!A$ هستند (یعنی
    می‌توانیم فرض~$!A$ را با برچسب~$n$ در~$\delta$ !!{discharge}).
  \end{enumerate}
\end{prop}

\begin{proof}
  باید نشان دهیم روابط متناظر میان اعداد گودل !!{formula}s و اعداد
  گودل !!{derivation}s بازگشتی اولیه‌اند.
  \begin{enumerate}
  \item می‌خواهیم نشان دهیم $\fn{Assum}(x,d,n)$، که وقتی $x$ عدد
    گودل فرضی با برچسب~$n$ در !!{derivation} دارای عدد گودل~$d$ باشد
    برقرار است، بازگشتی اولیه است. این درست است اگر !!{derivation} دارای
    عدد گودل~$\tuple{0,x,n}$ زیر!!{derivation}ی از~$d$ باشد. توجه کنید
    روش رمزگذاری !!{derivation}s حالت ویژه‌ای از رمزگذاری درخت‌ها در
    \olref[cmp][rec][tre]{sec} است؛ پس تابع بازگشتی اولیهٔ
    $\fn{SubtreeSeq}(d)$ دنباله‌ای از اعداد گودل همهٔ
    زیر!!{derivation}sی~$d$ (به طول حداکثر~$d$) به دست می‌دهد. بنابراین
    می‌توانیم تعریف کنیم
    \[
    \fn{Assum}(x, d, n) \defiff \bexists{i<d}{(\fn{SubtreeSeq}(d))_i =
      \tuple{0, x, n}}.
    \]
  \item می‌خواهیم نشان دهیم $\fn{Discharge}(x,d,n)$، که وقتی همهٔ
    فرض‌های دارای برچسب~$n$ در !!{derivation} دارای عدد گودل~$d$ همان
    !!{formula} دارای عدد گودل~$x$ باشند برقرار است. اما این رابطه
    برقرار است اگر و تنها اگر
    $\bforall{y<d}{(\fn{Assum}(y,d,n)\lif y=x)}$.
  \end{enumerate}
\end{proof}

\begin{prop}
  \ollabel{prop:followsby}
  ویژگی~$\fn{Correct}(d)$ که برقرار است اگر و تنها اگر آخرین استنتاج
  در !!{derivation}~$\delta$ با عدد گودل~$d$ درست باشد، بازگشتی اولیه است.
\end{prop}

\begin{proof}
  اینجا باید نشان دهیم به ازای هر قاعدهٔ استنتاج~$R$، رابطهٔ
  $\fn{FollowsBy}_R(d)$ بازگشتی اولیه است؛
  $\fn{FollowsBy}_R(d)$ برقرار است اگر و
  تنها اگر $d$ عدد گودل !!{derivation}~$\delta$ باشد و !!{formula} پایانی
  $\delta$ با کاربرد درست~$R$ از زیر!!{derivation}sی بی‌واسطهٔ~$\delta$
  نتیجه شود.

  حالت ساده قاعدهٔ~$\Intro{\land}$ است. اگر $\delta$ با استنتاج درست
  $\Intro{\land}$ پایان یابد چنین شکلی دارد:
  \begin{prooftree}
    \AxiomC{}
    \RightLabel{$\delta_1$}
    \DeduceC{$!A$}

    \AxiomC{}
    \RightLabel{$\delta_2$}
    \DeduceC{$!B$}

    \RightLabel{\Intro\land}
    \BinaryInfC{$!A \land !B$}
  \end{prooftree}
  آنگاه عدد گودل~$d$ برای~$\delta$ برابر
  $\tuple{2,d_1,d_2,\Gn{(!A\land !B)},0,k}$ است که در آن
  $\fn{EndFmla}(d_1)=\Gn{!A}$، $\fn{EndFmla}(d_2)=\Gn{!B}$، $n=0$
  و~$k=1$. پس می‌توانیم $\fn{FollowsBy}_{\Intro\land}(d)$ را چنین
  تعریف کنیم:
  \begin{multline*}
    (d)_0 = 2 \land \fn{DischargeLabel}(d) = 0 \land \fn{LastRule}(d) = 1 \land {}\\
    \fn{EndFmla}(d) = {}\\
    \Gn{(} \concat \fn{EndFmla}((d)_1) \concat \Gn{\land}
    \concat \fn{EndFmla}((d)_2) \concat \Gn{)}.
  \end{multline*}

  نمونهٔ سادهٔ دیگر قاعدهٔ~$\Intro\eq$ است. این قاعده مقدمه‌ای ندارد،
  پس مانند فرض‌ها $(d)_0=0$. همچنین برچسب رفعی ندارد؛ یعنی $n=0$.
  اما $!A$ باید برای ترم بسته‌ای چون~$t$ به صورت~$\eq[t][t]$ باشد.
  تعریف بازگشتی اولیه چنین است:
  \begin{multline*}
    (d)_0 = 0 \land \fn{DischargeLabel}(d) = 0 \land {}\\
    \bexists{t<d}{(\fn{ClTerm}(t) \land
      \fn{EndFmla}(d) = {}\\
      \Gn{{\eq}(} \concat t \concat \Gn{,} \concat t \concat \Gn{)})}.
  \end{multline*}

  برای نمونه‌ای پیچیده‌تر، $\fn{FollowsBy}_{\Intro{\lif}}(d)$ برقرار
  است اگر و تنها اگر !!{formula} پایانی~$\delta$ به صورت~$(!A\lif !B)$،
  !!{formula} پایانی~$\delta_1$ برابر~$!B$، و هر فرضِ~$\delta$ با
  برچسب~$n$ به صورت~$!A$ باشد. این را به‌طور بازگشتی اولیه چنین بیان
  می‌کنیم:
  \begin{multline*}
    (d)_0 = 1 \land {}\\
    \bexists{a<d}{(\fn{Discharge}(a, (d)_1, \fn{DischargeLabel}(d)) \land {}}\\
      \fn{EndFmla}(d) = (\Gn{(} \concat a \concat \Gn{\lif}
      \concat \fn{EndFmla}((d)_1) \concat \Gn{)}))
  \end{multline*}
  ($a$ را عدد گودل~$!A$ در نظر بگیرید).

  به‌عنوان نمونه‌ای دیگر، $\Intro{\lexists}$ را در نظر بگیرید. اینجا
  آخرین استنتاج در~$\delta$ درست است اگر و تنها اگر !!a{formula}~$!A$،
  ترم بستهٔ~$t$ و !!a{variable}~$x$ای وجود داشته باشند به‌گونه‌ای که
  $\Subst{!A}{t}{x}$ !!{formula} پایانی !!{derivation}~$\delta_1$ و
  $\lexists[x][!A]$ نتیجهٔ آخرین استنتاج باشد. بنابراین
  $\fn{FollowsBy}_{\Intro{\lexists}}(d)$ برقرار است اگر و تنها اگر
  \begin{multline*}
    (d)_0 = 1 \land \fn{DischargeLabel}(d) = 0 \land {} \\
    \bexists{a < d}{\bexists{x<d}{\bexists{t<d}{
          (\fn{ClTerm}(t) \land \fn{Var}(x)  \land {}}}}\\
    \fn{Subst}(a,t,x) = \fn{EndFmla}((d)_1) \land
    \fn{EndFmla}(d) = (\Gn{\lexists} \concat x \concat a)).
  \end{multline*}

  سپس $\fn{Correct}(d)$ را چنین تعریف می‌کنیم:
  \begin{multline*}
    \fn{Sent}(\fn{EndFmla}(d)) \land {}\\
    [ (\fn{LastRule}(d) = 1 \land
    \fn{FollowsBy}_{\Intro\land}(d)) \lor \dots \lor {}\\
    (\fn{LastRule}(d) = 16 \land \fn{FollowsBy}_{\Elim\eq}(d)) \lor {}\\
    \bexists{n<d}{\bexists{x<d}{(d = \tuple{0, x, n})}} ].
  \end{multline*}
  سطر نخست تضمین می‌کند !!{formula} پایانی~$d$ یک جمله است. سطر آخر
  حالتی را پوشش می‌دهد که $d$ فقط یک فرض است.
\end{proof}

\begin{prob}
  ویژگی‌های زیر را مانند \olref[inc][art][pnd]{prop:followsby} تعریف
  کنید:
  \begin{enumerate}
  \item $\fn{FollowsBy}_{\Elim{\lif}}(d)$؛
  \item $\fn{FollowsBy}_{\Elim{\eq}}(d)$؛
  \item $\fn{FollowsBy}_{\Elim{\lor}}(d)$؛
  \item $\fn{FollowsBy}_{\Intro{\lforall}}(d)$.
  \end{enumerate}
  برای مورد آخر باید همچنین نشان دهید می‌توان به‌طور بازگشتی اولیه
  آزمود که آخرین استنتاج !!{derivation} دارای عدد گودل~$d$ شرط متغیر ویژه
  را برآورده می‌کند؛ یعنی متغیر ویژهٔ~$a$ در استنتاج
  $\Intro{\lforall}$ نه در !!{formula} پایانی~$d$ و نه در فرض بازی از
  $d$ رخ می‌دهد. می‌توانید از محمول بازگشتی اولیهٔ~$\fn{OpenAssum}$
  در \olref[inc][art][pnd]{prop:openassum} استفاده کنید.
\end{prob}

\begin{prop}
  \ollabel{prop:deriv}
  رابطهٔ~$\fn{Deriv}(d)$ که وقتی $d$ عدد گودل !!{derivation} درست
  $\delta$ باشد برقرار است، بازگشتی اولیه است.
\end{prop}

\begin{proof}
  !!^a{derivation}~$\delta$ درست است اگر هر استنتاج آن کاربردی درست از یک
  قاعده باشد؛ یعنی هر زیر!!{derivation}s آن با استنتاجی درست پایان یابد.
  پس $\fn{Deriv}(d)$ اگر و تنها اگر
  \[
  \bforall{i<\len{\fn{SubtreeSeq}(d)}}{\fn{Correct}((\fn{SubtreeSeq}(d))_i)}.
  \]
\end{proof}

\begin{prop}
  \ollabel{prop:openassum} رابطهٔ~$\fn{OpenAssum}(z,d)$ که وقتی $z$
  عدد گودل یک فرض !!a{undischarged}ٔ~$!A$ در !!{derivation}~$\delta$ با عدد
  گودل~$d$ باشد برقرار است، بازگشتی اولیه است.
\end{prop}

\begin{proof}
  یک رخداد فرض هنگامی !!{discharged} است که با برچسب~$n$ در
  زیر!!{derivation}ی از~$\delta$ رخ دهد که با قاعده‌ای دارای برچسب رفع~$n$
  پایان می‌یابد. پس $!A$ فرضی !!a{undischarged} از~$\delta$ است اگر دست‌کم
  یکی از رخدادهایش در~$\delta$ !!{discharged} نباشد. باید مراقب باشیم:
  ممکن است $\delta$ هم رخدادهای !!{discharged} و هم رخدادهای
  !!{undischarged} از~$!A$ داشته باشد.

  دنبالهٔ~$\delta_0$،~\dots، $\delta_k$ را در نظر بگیرید که
  $\delta_0=\delta$، $\delta_k$ فرض~$\Discharge{!A}{n}$ (برای بعضی
  $n$)، و~$\delta_{i+1}$ زیر!!{derivation} بی‌واسطه‌ای از~$\delta_i$ است.
  اگر چنین دنباله‌ای وجود داشته باشد که در آن هیچ~$\delta_i$ با
  استنتاجی دارای برچسب رفع~$n$ پایان نیابد، آنگاه $!A$ فرضی
  !!a{undischarged} از~$\delta$ است.

  تابع بازگشتی اولیهٔ~$\fn{SubtreeSeq}(d)$ دنباله‌ای از اعداد گودل
  همهٔ زیر!!{derivation}sی~$\delta$ فراهم می‌کند. هر دنبالهٔ اعداد گودل
  زیر!!{derivation}sی~$\delta$ زیر‌دنباله‌ای از آن است. زیر‌دنباله‌بودن
  رابطه‌ای بازگشتی اولیه است: $\fn{Subseq}(s,s')$ برقرار است اگر و
  تنها اگر
  $\bforall{i<\len{s}}{\lexists[j<\len{s'}][(s)_i=(s')_j]}$.
  زیر!!{derivation} بی‌واسطه بودن نیز چنین است:
  $\fn{Subderiv}(d,d')$ اگر و تنها اگر
  $\bexists{j<(d')_0}{d=(d')_{j+1}}$. پس می‌توانیم
  $\fn{OpenAssum}(z,d)$ را چنین تعریف کنیم:
  \begin{multline*}
    \bexists{s\leq\fn{SubtreeSeq}(d)}{(\fn{Subseq}(s, \fn{SubtreeSeq}(d))
      \land (s)_0 = d \land {}} \\
    \bexists{n<d}{((s)_{\len{s} \tsub 1} = \tuple{0, z, n} \land {}}\\
      \bforall{i<(\len{s} \tsub 1)}{(\fn{Subderiv}((s)_{i+1}, (s)_i) \land {}}\\
      \fn{DischargeLabel}((s)_i) \neq n))).
  \end{multline*}
\end{proof}

\begin{prop}
  \ollabel{prop:prf-prim-rec}
  فرض کنید $\Gamma$ مجموعه‌ای بازگشتی اولیه از !!{sentence}s باشد. آنگاه
  رابطهٔ~$\Prf[\Gamma](x,y)$ که بیان می‌کند «$x$ رمز
  !!a{derivation}~$\delta$ از~$!A$ بر پایهٔ فرض‌های !!{undischarged} در
  $\Gamma$ است و~$y$ عدد گودل~$!A$ است»، بازگشتی اولیه است.
\end{prop}

\begin{proof}
  فرض کنید «$y\in\Gamma$» با محمول بازگشتی اولیهٔ~$R_\Gamma(y)$ داده
  شود. باید نشان دهیم $\Prf[\Gamma](x,y)$، که برقرار است اگر و تنها
  اگر $y$ عدد گودل جملهٔ~$!A$ و~$x$ رمز !!{derivation}ی در استنتاج طبیعی
  با !!{formula} پایانی~$!A$ و همهٔ فرض‌های !!{undischarged} در~$\Gamma$
  باشد، بازگشتی اولیه است.

  بنا بر \olref{prop:deriv}، ویژگی~$\fn{Deriv}(x)$ که برقرار است اگر
  و تنها اگر $x$ عدد گودل !!{derivation} درست~$\delta$ در استنتاج طبیعی
  باشد بازگشتی اولیه است. پس می‌توانیم $\Prf[\Gamma](x,y)$ را چنین
  تعریف کنیم:
  \begin{align*}
    \Prf[\Gamma](x, y) \defiff {}
    & \fn{Deriv}(x) \land \fn{EndFmla}(x) = y \land {} \\
    & \bforall{z < x}{(\fn{OpenAssum}(z, x) \lif R_\Gamma(z))}.
  \end{align*}
\end{proof}

\end{document}
